%%%%%%%%%%%%%%%%%%%%%%%%%%%%%%%%%%%%%%%%%
% Developer CV
%----------------------------------------------------------------------------------------
%	PACKAGES AND OTHER DOCUMENT CONFIGURATIONS
%----------------------------------------------------------------------------------------

\documentclass[9pt]{developercv} % Default font size, values from 8-12pt are recommended
\usepackage{multicol}
\setlength{\columnsep}{0mm}
%----------------------------------------------------------------------------------------
\usepackage{lipsum}  


\begin{document}

%----------------------------------------------------------------------------------------
%	TITLE AND CONTACT INFORMATION
%----------------------------------------------------------------------------------------

\begin{minipage}[t]{0.5\textwidth} 
	\vspace{-\baselineskip} % Required for vertically aligning minipages
	
	{ \fontsize{16}{20} \textcolor{black}{\textbf{\MakeUppercase{Andrey Otvagin}}}} % First name
	
	\vspace{6pt}
    \icon{Envelope}{11}{\href{mailto:otvagin@ucsb.edu}{otvagin@ucsb.edu}}
    \hspace{12pt}

    \icon{Github}{11}{\href{https://github.com/an0tv}{github.com/an0tv}}
    \hspace{12pt}
    
    \icon{Chain}{11}{\href{https://an0tv.github.io/}{an0tv.github.io}}
    \hspace{12pt}

    
\end{minipage}
\hfill
\begin{minipage}[t]{0.5\textwidth}
    \vspace{-18pt}
    \cvsect{Skills}
    \vspace{-6pt}
    
    \begin{minipage}[t]{0.2\textwidth}
        \textbf{Languages:}
    \end{minipage}
    \hfill
    \begin{minipage}[t]{0.73\textwidth}
        C, C++, Python, ASM, SystemVerilog, VHDL, R
    \end{minipage}
    \vspace{4mm}
    
    \begin{minipage}[t]{0.2\textwidth}
        \textbf{Technologies:}
    \end{minipage}
    \hfill
    \begin{minipage}[t]{0.73\textwidth}
        Docker, Git, Eclipse Mosquitto, MQTT, WSL, OpenOCD, KiCad, Altium, ModelSim, LTspice, Logisim-Evolution, Vivado, Cadence Virtuoso, AWS, Digital Ocean
    \end{minipage}
    
\end{minipage}


%----------------------------------------------------------------------------------------
%	EDUCATION
%----------------------------------------------------------------------------------------
\vspace{-8 pt}
\cvsect{Education}
\begin{entrylist}

\vspace{-4pt}

    \entry
    {}
    {University of California, Santa Barbara}
    {September 2025 – June 2027}
    {M.S. Computer Engineer}
    \entry
    {}
    {University of California, Santa Barbara}
    {September 2022 – June 2026}
    {B.S. Computer Engineering and B.S. Statistics \& Data Science 
    \begin{itemize}[noitemsep,topsep=0pt,parsep=0pt,partopsep=0pt, leftmargin=10pt]
        \item Relevant Coursework (Computer Engineering): Deep Learning, Advanced Computer Architecture, Data Structures \& Algorithms, Embedded System, Machine Learning, Circuits, VLSI
        \item Relevant Coursework (Statistics \& Data Science): SAS Programming, SQL, Advanced Statistical Models, Survival Analysis, Stochastic Processes, Statistical Machine Learning, Regression Analysis, Advanced Linear Algebra
    \end{itemize}}
    
    
    
\end{entrylist}

%----------------------------------------------------------------------------------------
%	EXPERIENCE
%----------------------------------------------------------------------------------------
\vspace{-16 pt}
\cvsect{Experience}
\begin{entrylist}
    \entry
        {}
        {Becton Dickinson Bioscience - Firmware R\&D Engineer}
        {June 2024 – August 2025}
        {\vspace{-8pt}
        \begin{itemize}[noitemsep,topsep=0pt,parsep=0pt,partopsep=0pt, leftmargin=10pt]
            \item Developed Hardware and Software for HOOTL and HITL test suite for embedded systems. 
            \item Optimized system responsiveness by designing timeout strategies for long-running processes, enhancing overall system reliability. 
            \item Deployed a build pipeline to allow for nightly firmware builds for embedded subsystems.  
            \item Created thread-safe asynchronous fault monitoring system for embedded linux environment. 
        \end{itemize}}
    \entry
        {}
        {Gaucho Racing (UCSB Formula SAE) - Electronics \& Data Lead}
        {September 2022 – Present}
        {\vspace{-8pt}
        \begin{itemize}[noitemsep,topsep=0pt,parsep=0pt,partopsep=0pt, leftmargin=10pt]
            \item Led development of SAM, a high-throughput sensor array module for automotive systems, enabling versatile data collection including tire and rotor temperature, push-rod force, ride height, and vehicle dynamics monitoring via EKF with differential GPS and accelerometers.
            \item Deployed a real-time data collection \& analytics platform on Jetson Orin Nano to monitor data streamed by SAM through a custom CAN \& CANFD protocol translator, allowing for real-time thermal prediction and anomaly detection. Data collection and analytics platform built using MQTT, AWS, and SingleStore. 
            \item Engineered a custom CAN protocol with built-in system prioritization, ensuring reliable communication under high bus load. 
            \item Implemented CI/CD pipelines for a HOOTL system, accelerating development and reducing integration friction.
            \item Promoted collaborative coding practices and knowledge sharing within the team through
            extensive documentation.
            \item Coordinated closely with the mechanical and aerodynamic teams for better hardware integration into respective assemblies.  
            
        \end{itemize}}
\end{entrylist}

%----------------------------------------------------------------------------------------
%	Projects
%----------------------------------------------------------------------------------------
\vspace{-16 pt}
\cvsect{Projects}
\begin{entrylist}
    \entry
    {}
    {CHIRP - Chirp radar fusion and tracking }
    {Fall 2025 - Present}
    {\vspace{-4pt}
    \begin{itemize}[noitemsep,topsep=0pt,parsep=0pt,partopsep=0pt, leftmargin=10pt]
        \item Scalable chirp radar fusion architecture to support accurate multi-object tracking across multiple radar sensors, based on DCA1000EVM and AWR2243Boost hardware stack, with Jetson Orin Nano used to process sensor data. 
        \item Current features include multi-node synced CFAR, DBSCAN, and EKF pipeline running at 15FPS. 
    \end{itemize}
    }
    
    \entry
    {}
    {RISC Processor}
    {Spring 2025}
    {\vspace{-4pt}
    \begin{itemize}[noitemsep,topsep=0pt,parsep=0pt,partopsep=0pt, leftmargin=10pt]
    \item Created a pipelined RISC-V processor in verilog.
    \item Features include: Hazard Detection with respective forwarding and flushing, Branch Prediction implemented with  5-bit GHR, 2-bit PHT (32 entries), and 32-entry BTB, Cache emulation with Hardware Prefetcher as well as Critical Word First and Early Restart, and Superscalar implemented via doubling the stages in the pipeline and modifying hazard and forwarding to cover cross pipeline forwarding as well as single branch execution.
    \end{itemize}
    }

    \entry
    {}
    {SAM}
    {Fall 2024 - Winter 2026}
    {\vspace{-8pt}
    \begin{itemize}[noitemsep,topsep=0pt,parsep=0pt,partopsep=0pt, leftmargin=10pt]
    \item Teensy 4.1 alternative using an STM32G4 with a focus on automotive usage.
    \item Designed a compatible set of sub-modules for SAE vehicle characterization, meant for dynamic testing centered around small design and easy implementation. 
    \end{itemize}
    }
\end{entrylist}
\vspace{-16 pt}
%----------------------------------------------------------------------------------------
%	Publications
%----------------------------------------------------------------------------------------
\cvsect{Publications}
\begin{entrylist}
    \entry
		{}
		{Mapache - An Intelligent, Real-Time Telemetry Platform for Formula SAE}
		{}
		{\vspace{-8pt}
        \begin{itemize}[noitemsep,topsep=0pt,parsep=0pt,partopsep=0pt, leftmargin=10pt]
            \item Bharat Kathi, Andrey Otvagin, Jacob Jurek, Austin Chan, Colin Bickel
            \item Accepted to the 2025 International Telemetering Conference (ITC), to be published in their Vol. 60 Proceedings
            \item https://github.com/Gaucho-Racing/Mapache/blob/main/itc25.pdf
        \end{itemize}}
\end{entrylist}
\end{document}
